\chapter{Introdução}

\setlength{\epigraphwidth}{0.6\textwidth}

\epigraph{\justifying Quatro homens, cegos de nascença, queriam saber como era um elefante; eles pediram, então, informações a um tratador de elefantes. Ele os levou até um elefante e os convidou a examiná-lo; um dos homens apalpou a pata do elefante, outro a tromba, outro a cauda, e o quarto a orelha. Então, tentaram descrever o elefante uns aos outros. O primeiro homem disse: `O elefante é como uma árvore'. `Não', disse o segundo, `o elefante é como uma cobra'. `Bobagem!', disse o terceiro, `o elefante é como uma vassoura'. `Vocês estão todos errados', disse o quarto, `o elefante é como um leque'. E assim continuaram discutindo entre si, enquanto o elefante os observava em silêncio.}{\textcite[p. vii]{Johnstone2002}}


Há uma velha parábola hindu sobre sábios cegos que, um dia, tiveram a vontade de presenciar um elefante. Quando apresentados ao animal, cada um foi para um lado e o apreendeu apalpando-o. Pelo que os seus sentidos lhes forneciam, eles concluíram, cada um em sua particularidade, que o animal era tal como ele se lhes apresentava. Ao averiguar que suas apreensões não coincidiam entre si, tinham de concluir que os demais estavam incorretos, porquanto uma das poucas verdades de que tinham certeza era aquilo que seus sentidos lhes permitiam aferir da realidade. 

Embora a lição de moral desta parábola se trate, originalmente, da unicidade de Deus e da futilidade de sectarismos quanto à descrição de sua natureza,\footnote{``É assim que sectários, influenciados por seus respectivos sistemas, discutem sobre a natureza de Deus, a qual a mente não consegue alcançar'' \parencite[p. 30]{Robinson1885}; ``Em tais guerras teológicas // é comum, receio eu crer, // haver discussões espúrias // do que o outro quer dizer: // tagarelam do Elefante // o qual nenhum pôde ver!'' \parencite[tradução minha do poema]{Saxe1872}. (Todas as citações traduzidas neste trabalho, cujas obras citadas não estejam em português, são de minha autoria, a menos de menção explícita.)} ela ainda serve como uma história de advertência [\textit{cautionary tale}] sobre o fazer científico, em particular quando seu objeto de estudo é tão complexo quanto a economia. Sua principal conclusão, porém, é de suma importância para a crítica a ser feita no presente trabalho: os complexos parciais de uma totalidade podem assumir funções que sejam opostas umas às outras, de tal forma que sua compreensão ``genuína'' advém necessariamente de sua intelecção \textit{enquanto partes desta totalidade}. 

De fato, o propósito daquela parábola é de que remetamos à imagem, que já temos em mente, do elefante enquanto um só ente, a partir da qual \textit{sabemos} que as visões dos sábios cegos são não mais que parciais; conseguimos ver como e por que suas visões são parciais, que um está apreendendo ``uma pata'', ``uma tromba'' etc., justamente porque é ``uma pata'', ``uma tromba'' etc. \textit{deste ente-elefante}. A conclusão dessa estória é que, embora suas apreensões particulares não estejam totalmente erradas, elas somente assumem um caráter razoável de verdade \textit{quando são remetidas ao todo ao qual pertencem}. Nas Ciências Exatas, isso se resume à afirmação de que ``o caráter local/particular de um objeto não tem de coincidir com seu caráter global/universal'': o exemplo canônico é que a Terra é \textit{localmente} plana, mas é \textit{globalmente} esferoide, fato o qual permite entender \textit{por que} ela é localmente plana.\footnote{Esta explicação é rigorosa: matematicamente, uma esfera, vista como uma variedade diferenciável, possui dimensão $2$, o que quer dizer que qualquer região (aberta) em sua superfície é homeomorfa (``idêntica em forma'') a um plano bidimensional em $\mathbb{R}^2$. É por este motivo que conseguimos mapear a Terra, a princípio algo tridimensional, através de vários mapas e atlas bidimensionais. Por outro lado, é justamente devido à sua geometria esférica que não é possível mapeá-la através de \textit{um único} mapa-mundi (bidimensional) sem haver distorções: porque uma superfície plana e uma superfície curva possuem \textit{curvaturas} (também uma quantidade matematicamente bem-definida) diferentes, o que impede que elas sejam \textit{globalmente} homeomorfas. (Diga-se de passagem que a presente discussão resvala, de tabela, nas teorias da conspiração sobre a Terra plana: um terraplanista não está totalmente incorreto quando diz que ``a Terra é plana'', posto que é assim que ela \textit{se apresenta aos nossos sentidos}, mas erra quando extrapola esta conclusão empírica à Terra \textit{em sua totalidade}.)}

No que tange à descrição da economia capitalista, é esclarecedor por que Marx primeiramente analisa os ciclos do capital-dinheiro,\footnote{Ao longo deste trabalho, usar-se-á a tradução ``capital-dinheiro'' para a categoria \textit{Geldkapital} invés de ``capital monetário'', como presente nas edições da Boitempo d'\textit{O Capital} \parencite{Marx2014,Marx2017a}. Essa tradução evita a confusão de que esta categoria esteja associada à \textit{moeda} (``monetário''), pois trata-se somente da forma-\textit{dinheiro} --- equivalente universal de valor --- que o capital assume em seu movimento próprio. Devo esta correção terminológica à certeira insistência do prof. Marcelo Carcanholo em suas aulas.} capital produtivo e capital-mercadoria, para só depois apresentar o conceito de capital industrial: ele afere a observação parcial deste movimento total, também através da visão dos economistas políticos do passado, reconhece a veracidade destas formas, \textit{mas não se restringe a elas}, buscando, portanto, o ``elefante'' ao qual estas formas remetem. Fica claro aqui, inclusive, em que sentido exatamente Marx fala sobre como as ciências sociais não podem se valer de experimentos e somente da ``força da abstração'': não se pode apreender \textit{empiricamente} o movimento do capital industrial em sua inteireza, somente suas formas fenomenais e, por sua própria natureza, \textit{parciais}, sendo-nos necessário o ``olho da razão'' para entrever o todo ao qual tais formas remetem, e das quais emanam. 

